\section{Exercices d'application}

\rubrique{Barycentre de deux points pondérés}

\exo{}\diff{1}
Soient $A$ et $B$ deux points du plan. Construire le barycentre $G$ des points pondérés $(A,3)$ et $(B,-1)$.

\exo{}\diff{1}
Déterminer les réels $\alpha$ et $\beta$ tels que $G$ soit le barycentre de $(A,\alpha)$ et $(B,\beta)$ sachant que $\overrightarrow{AG} = \dfrac{2}{5}\overrightarrow{AB}$.

\exo{}\diff{2}
On donne $A(2,3)$ et $B(-1,4)$ dans un repère orthonormé $(O,\vec{i},\vec{j})$. Calculer les coordonnées du barycentre $G$ de $(A,2)$ et $(B,3)$.

\rubrique{Barycentre de $n$ points pondérés — Associativité}

\exo{}\diff{1}
Soient $A$, $B$, $C$ trois points du plan. Construire le barycentre $G$ de $(A,1)$, $(B,2)$ et $(C,3)$.

\exo{}\diff{2}
On considère les points $A$, $B$, $C$ et $D$. Soit $G$ le barycentre de $(A,2)$, $(B,1)$, $(C,-1)$ et $(D,3)$. En utilisant l'associativité, déterminer le barycentre partiel de $(A,2)$ et $(B,1)$ ; en déduire une relation vectorielle simple.

\exo{}\diff{2}
$ABC$ est un triangle. Soit $G$ le barycentre de $(A,2)$, $(B,3)$ et $(C,5)$. Montrer que $G$ appartient à la médiane issue de $C$.

\rubrique{Réduction de sommes vectorielles}

\exo{}\diff{2}
Soient $A$, $B$, $C$ trois points du plan. Réduire la somme vectorielle $\overrightarrow{MA} + 2\overrightarrow{MB} - 3\overrightarrow{MC}$ en utilisant le barycentre $G$ des points $(A,1)$, $(B,2)$ et $(C,-3)$.

\exo{}\diff{2}
Exprimer $3\overrightarrow{MA} - 2\overrightarrow{MB} + \overrightarrow{MC}$ sous la forme $k\overrightarrow{MG}$, où $G$ est un point fixe à déterminer.

\exo{}\diff{3}
Soient $A$, $B$, $C$ trois points distincts. Déterminer l'ensemble des points $M$ du plan tels que $\|\overrightarrow{MA} + 2\overrightarrow{MB} - 3\overrightarrow{MC}\| = 4$.

\rubrique{Lignes de niveau : $M \mapsto MA^2 + MB^2$}

\exo{}\diff{2}
Soient $A$ et $B$ deux points distincts tels que $AB = 6$. Déterminer l'ensemble des points $M$ du plan vérifiant $MA^2 + MB^2 = 50$.

\exo{}\diff{2}
On donne $A(1,2)$ et $B(5,4)$. Déterminer l'ensemble des points $M$ tels que $MA^2 + MB^2 = 40$.

\exo{}\diff{3}
$ABC$ est un triangle rectangle en $A$ avec $AB = 4$ et $AC = 3$. Déterminer l'ensemble des points $M$ du plan tels que $MA^2 + MB^2 + MC^2 = 50$.

\rubrique{Lignes de niveau : $M \mapsto \dfrac{MA}{MB}$}

\exo{}\diff{2}
Soient $A$ et $B$ deux points distincts. Déterminer l'ensemble des points $M$ du plan tels que $\dfrac{MA}{MB} = 2$.

\exo{}\diff{2}
On donne $A(-1,0)$ et $B(3,0)$. Déterminer l'ensemble des points $M$ tels que $MA = 3MB$.

\exo{}\diff{3}
$ABC$ est un triangle. Déterminer l'ensemble des points $M$ du plan tels que $\dfrac{MA}{MB} = \dfrac{MC}{MA}$.

\section{Exercices d'entraînement}

\rubrique{Barycentre de deux points pondérés}

\exo{}\diff{1}
Construire le barycentre $G$ de $(A,2)$ et $(B,3)$.

\exo{}\diff{1}
Déterminer les coefficients $a$ et $b$ tels que $G$ soit le barycentre de $(A,a)$ et $(B,b)$ sachant que $\overrightarrow{BG} = 3\overrightarrow{BA}$.

\exo{}\diff{2}
Dans un repère, on donne $A(-2,1)$ et $B(4,3)$. Calculer les coordonnées du barycentre $G$ de $(A,5)$ et $(B,-2)$.

\exo{}\diff{2}
Soient $A$ et $B$ deux points. Déterminer l'ensemble des points $M$ tels que $\overrightarrow{MA} + 2\overrightarrow{MB} = \vec{0}$.

\rubrique{Barycentre de $n$ points pondérés}

\exo{}\diff{1}
Construire le barycentre $G$ de $(A,1)$, $(B,1)$ et $(C,1)$ (isobarycentre).

\exo{}\diff{2}
Soient $A$, $B$, $C$, $D$ quatre points. Soit $G$ le barycentre de $(A,2)$, $(B,3)$, $(C,-1)$ et $(D,4)$. Déterminer le barycentre partiel de $(C,-1)$ et $(D,4)$ ; en déduire une relation.

\exo{}\diff{2}
$ABC$ est un triangle. Soit $G$ le barycentre de $(A,1)$, $(B,2)$ et $(C,3)$. Montrer que $G$ est le point d'intersection des médianes.

\exo{}\diff{3}
On considère un quadrilatère $ABCD$. Soit $G$ le barycentre de $(A,1)$, $(B,2)$, $(C,3)$ et $(D,4)$. Montrer que $G$ est le centre de gravité du triangle formé par les barycentres partiels.

\rubrique{Réduction de sommes vectorielles}

\exo{}\diff{2}
Réduire $\overrightarrow{MA} - 2\overrightarrow{MB} + 3\overrightarrow{MC}$.

\exo{}\diff{2}
Exprimer $2\overrightarrow{MA} + 3\overrightarrow{MB} - \overrightarrow{MC}$ sous la forme $k\overrightarrow{MG}$.

\exo{}\diff{3}
Déterminer l'ensemble des points $M$ tels que $\|2\overrightarrow{MA} - \overrightarrow{MB} + \overrightarrow{MC}\| = 5$.

\exo{}\diff{3}
Soient $A$, $B$, $C$ trois points. Déterminer l'ensemble des points $M$ tels que $\overrightarrow{MA} + 2\overrightarrow{MB} + 3\overrightarrow{MC} = \vec{0}$.

\rubrique{Lignes de niveau : $M \mapsto MA^2 + MB^2$}

\exo{}\diff{2}
Soient $A$ et $B$ tels que $AB = 8$. Déterminer l'ensemble des points $M$ tels que $MA^2 + MB^2 = 100$.

\exo{}\diff{2}
On donne $A(0,0)$ et $B(6,0)$. Déterminer l'ensemble des points $M$ tels que $MA^2 + MB^2 = 72$.

\exo{}\diff{3}
$ABC$ est un triangle équilatéral de côté $4$. Déterminer l'ensemble des points $M$ tels que $MA^2 + MB^2 + MC^2 = 48$.

\exo{}\diff{3}
Soient $A$, $B$, $C$ trois points. Déterminer l'ensemble des points $M$ tels que $MA^2 + MB^2 = MC^2$.

\rubrique{Lignes de niveau : $M \mapsto \dfrac{MA}{MB}$}

\exo{}\diff{2}
Déterminer l'ensemble des points $M$ tels que $\dfrac{MA}{MB} = 3$.

\exo{}\diff{2}
On donne $A(2,0)$ et $B(8,0)$. Déterminer l'ensemble des points $M$ tels que $MA = 2MB$.

\exo{}\diff{3}
$ABC$ est un triangle. Déterminer l'ensemble des points $M$ tels que $\dfrac{MA}{MB} = \dfrac{MB}{MC}$.

\exo{}\diff{3}
Soient $A$ et $B$ deux points. Déterminer l'ensemble des points $M$ tels que $\dfrac{MA}{MB} = k$, où $k$ est un réel positif différent de $1$.

\section{Problèmes type Probatoire/Bac}

\sujetbac[Probatoire Blanc — Barycentre et lignes de niveau]
\label{prob:1}
Dans le plan muni d'un repère orthonormé $(O,\vec{i},\vec{j})$, on considère les points $A(2,1)$, $B(4,5)$ et $C(-2,3)$.

\begin{enumerate}[label=\arabic*)]
\item Calculer les coordonnées du barycentre $G$ de $(A,1)$, $(B,2)$ et $(C,3)$.
\item Soit $M$ un point quelconque du plan. Réduire la somme vectorielle $\overrightarrow{MA} + 2\overrightarrow{MB} + 3\overrightarrow{MC}$.
\item Déterminer l'ensemble des points $M$ du plan tels que $\|\overrightarrow{MA} + 2\overrightarrow{MB} + 3\overrightarrow{MC}\| = 6\sqrt{2}$.
\item Déterminer l'ensemble des points $M$ du plan tels que $MA^2 + 2MB^2 + 3MC^2 = 100$.
\end{enumerate}

\sujetbac[Baccalauréat Blanc — Lignes de niveau]
\label{prob:2}
Soient $A$ et $B$ deux points distincts du plan tels que $AB = 4$.

\begin{enumerate}[label=\arabic*)]
\item Déterminer l'ensemble $E_1$ des points $M$ du plan tels que $MA^2 + MB^2 = 20$.
\item Déterminer l'ensemble $E_2$ des points $M$ du plan tels que $\dfrac{MA}{MB} = 2$.
\item Soit $C$ un point tel que $ABC$ soit un triangle rectangle en $A$ avec $AC = 3$. Déterminer l'ensemble $E_3$ des points $M$ du plan tels que $MA^2 + MB^2 + MC^2 = 50$.
\item Montrer que les ensembles $E_1$, $E_2$ et $E_3$ sont concourants.
\end{enumerate}

\sujetbac[Probatoire — Barycentre et alignement]
\label{prob:3}
$ABC$ est un triangle. On considère les points $I$, $J$ et $K$ définis par :
\[
\overrightarrow{AI} = \dfrac{1}{3}\overrightarrow{AB}, \quad \overrightarrow{BJ} = \dfrac{2}{3}\overrightarrow{BC}, \quad \overrightarrow{CK} = \dfrac{3}{4}\overrightarrow{CA}.
\]

\begin{enumerate}[label=\arabic*)]
\item Exprimer $I$ comme barycentre de $A$ et $B$.
\item Exprimer $J$ comme barycentre de $B$ et $C$.
\item Exprimer $K$ comme barycentre de $C$ et $A$.
\item Soit $G$ le barycentre de $(A,2)$, $(B,1)$ et $(C,3)$. Montrer que $G$ est le point d'intersection des droites $(AI)$ et $(BJ)$.
\item Déterminer l'ensemble des points $M$ du plan tels que $2MA^2 + MB^2 + 3MC^2 = 2GA^2 + GB^2 + 3GC^2$.
\end{enumerate}

\section{Corrigés}

\corrige{de l'exercice 1}
Soit $G$ le barycentre de $(A,3)$ et $(B,-1)$. On a $3\overrightarrow{GA} - \overrightarrow{GB} = \vec{0}$, soit $3\overrightarrow{GA} - (\overrightarrow{GA} + \overrightarrow{AB}) = \vec{0}$, d'où $2\overrightarrow{GA} - \overrightarrow{AB} = \vec{0}$, donc $\overrightarrow{AG} = \dfrac{1}{2}\overrightarrow{AB}$. $G$ est le milieu de $[AB]$.

\corrige{de l'exercice 2}
On a $\overrightarrow{AG} = \dfrac{2}{5}\overrightarrow{AB}$, donc $G$ est le barycentre de $(A,3)$ et $(B,2)$ car $\overrightarrow{AG} = \dfrac{2}{5}\overrightarrow{AB} \iff 5\overrightarrow{AG} = 2\overrightarrow{AB} \iff 5\overrightarrow{AG} = 2(\overrightarrow{AG} + \overrightarrow{GB}) \iff 3\overrightarrow{GA} + 2\overrightarrow{GB} = \vec{0}$. Donc $\alpha = 3$, $\beta = 2$.

\corrige{de l'exercice 3}
$G$ barycentre de $(A,2)$ et $(B,3)$ donne $\overrightarrow{OG} = \dfrac{2\overrightarrow{OA} + 3\overrightarrow{OB}}{5}$. Ainsi $x_G = \dfrac{2\times 2 + 3\times (-1)}{5} = \dfrac{4 - 3}{5} = \dfrac{1}{5}$, $y_G = \dfrac{2\times 3 + 3\times 4}{5} = \dfrac{6 + 12}{5} = \dfrac{18}{5}$. Donc $G\left(\dfrac{1}{5},\dfrac{18}{5}\right)$.

\corrige{de l'exercice 4}
$G$ barycentre de $(A,1)$, $(B,2)$ et $(C,3)$. On a $\overrightarrow{GA} + 2\overrightarrow{GB} + 3\overrightarrow{GC} = \vec{0}$. En utilisant l'associativité, soit $H$ le barycentre de $(A,1)$ et $(B,2)$, alors $H$ est tel que $\overrightarrow{HA} + 2\overrightarrow{HB} = \vec{0}$, donc $\overrightarrow{AH} = \dfrac{2}{3}\overrightarrow{AB}$. Alors $G$ est le barycentre de $(H,3)$ et $(C,3)$, donc $G$ est le milieu de $[HC]$. Construction : placer $H$ sur $[AB]$ au $\dfrac{2}{3}$ à partir de $A$, puis $G$ milieu de $[HC]$.

\corrige{de l'exercice 5}
Soit $H$ le barycentre partiel de $(A,2)$ et $(B,1)$. Alors $H$ est tel que $2\overrightarrow{HA} + \overrightarrow{HB} = \vec{0}$, soit $\overrightarrow{AH} = \dfrac{1}{3}\overrightarrow{AB}$. $G$ est alors le barycentre de $(H,3)$ et $(C,-1)$ et $(D,3)$. On a $3\overrightarrow{GH} - \overrightarrow{GC} + 3\overrightarrow{GD} = \vec{0}$, soit $3\overrightarrow{GH} + 3\overrightarrow{GD} = \overrightarrow{GC}$, d'où $\overrightarrow{GC} = 3(\overrightarrow{GH} + \overrightarrow{GD})$.

\corrige{de l'exercice 6}
$G$ barycentre de $(A,2)$, $(B,3)$ et $(C,5)$. On a $2\overrightarrow{GA} + 3\overrightarrow{GB} + 5\overrightarrow{GC} = \vec{0}$. Soit $I$ le milieu de $[AB]$, alors $\overrightarrow{GA} + \overrightarrow{GB} = 2\overrightarrow{GI}$. Donc $2\overrightarrow{GA} + 3\overrightarrow{GB} = 2(\overrightarrow{GA} + \overrightarrow{GB}) + \overrightarrow{GB} = 4\overrightarrow{GI} + \overrightarrow{GB}$. L'équation devient $4\overrightarrow{GI} + \overrightarrow{GB} + 5\overrightarrow{GC} = \vec{0}$. En écrivant $\overrightarrow{GB} = \overrightarrow{GI} + \overrightarrow{IB}$, on a $5\overrightarrow{GI} + \overrightarrow{IB} + 5\overrightarrow{GC} = \vec{0}$, soit $5(\overrightarrow{GI} + \overrightarrow{GC}) + \overrightarrow{IB} = \vec{0}$. Soit $J$ le milieu de $[IC]$, alors $\overrightarrow{GI} + \overrightarrow{GC} = 2\overrightarrow{GJ}$, donc $10\overrightarrow{GJ} + \overrightarrow{IB} = \vec{0}$, soit $\overrightarrow{GJ} = -\dfrac{1}{10}\overrightarrow{IB}$. $G$ est sur la droite $(IC)$, qui est la médiane issue de $C$ car $I$ est le milieu de $[AB]$.

\corrige{de l'exercice 7}
Soit $G$ le barycentre de $(A,1)$, $(B,2)$ et $(C,-3)$. Alors $\overrightarrow{GA} + 2\overrightarrow{GB} - 3\overrightarrow{GC} = \vec{0}$. Pour tout point $M$, $\overrightarrow{MA} + 2\overrightarrow{MB} - 3\overrightarrow{MC} = (\overrightarrow{MG} + \overrightarrow{GA}) + 2(\overrightarrow{MG} + \overrightarrow{GB}) - 3(\overrightarrow{MG} + \overrightarrow{GC}) = (1+2-3)\overrightarrow{MG} + (\overrightarrow{GA} + 2\overrightarrow{GB} - 3\overrightarrow{GC}) = \vec{0}$. Donc la somme est nulle pour tout $M$.

\corrige{de l'exercice 8}
Soit $G$ le barycentre de $(A,3)$, $(B,-2)$ et $(C,1)$. Alors $3\overrightarrow{GA} - 2\overrightarrow{GB} + \overrightarrow{GC} = \vec{0}$. Pour tout $M$, $3\overrightarrow{MA} - 2\overrightarrow{MB} + \overrightarrow{MC} = (3-2+1)\overrightarrow{MG} = 2\overrightarrow{MG}$.

\corrige{de l'exercice 9}
Soit $G$ le barycentre de $(A,1)$, $(B,2)$ et $(C,-3)$. D'après l'exercice 7, $\overrightarrow{MA} + 2\overrightarrow{MB} - 3\overrightarrow{MC} = \vec{0}$ pour tout $M$, donc l'ensemble des points $M$ est le plan tout entier. Mais l'équation $\|\vec{0}\| = 4$ est impossible, donc l'ensemble est vide.

\corrige{de l'exercice 10}
Soit $I$ le milieu de $[AB]$. On a $MA^2 + MB^2 = 2MI^2 + \dfrac{AB^2}{2}$. Avec $AB = 6$, $MA^2 + MB^2 = 2MI^2 + 18$. L'équation $2MI^2 + 18 = 50$ donne $MI^2 = 16$, soit $MI = 4$. L'ensemble est le cercle de centre $I$ et de rayon $4$.

\corrige{de l'exercice 11}
$I$ milieu de $[AB]$ a pour coordonnées $(3,3)$. $AB^2 = (5-1)^2 + (4-2)^2 = 16 + 4 = 20$. $MA^2 + MB^2 = 2MI^2 + \dfrac{AB^2}{2} = 2MI^2 + 10$. L'équation $2MI^2 + 10 = 40$ donne $MI^2 = 15$, soit $MI = \sqrt{15}$. L'ensemble est le cercle de centre $I(3,3)$ et de rayon $\sqrt{15}$.

\corrige{de l'exercice 12}
$ABC$ rectangle en $A$, donc $BC^2 = AB^2 + AC^2 = 16 + 9 = 25$. Soit $G$ l'isobarycentre de $A$, $B$, $C$ (centre de gravité). On a $MA^2 + MB^2 + MC^2 = 3MG^2 + GA^2 + GB^2 + GC^2$. Or $GA^2 + GB^2 + GC^2 = \dfrac{AB^2 + AC^2 + BC^2}{3} = \dfrac{16 + 9 + 25}{3} = \dfrac{50}{3}$. Donc $3MG^2 + \dfrac{50}{3} = 50$, soit $3MG^2 = \dfrac{100}{3}$, $MG^2 = \dfrac{100}{9}$, $MG = \dfrac{10}{3}$. L'ensemble est le cercle de centre $G$ et de rayon $\dfrac{10}{3}$.

\corrige{de l'exercice 13}
$\dfrac{MA}{MB} = 2$ équivaut à $MA^2 = 4MB^2$, soit $MA^2 - 4MB^2 = 0$. Soit $I$ le barycentre de $(A,1)$ et $(B,-4)$ (car $1-4 \neq 0$). Alors $MA^2 - 4MB^2 = (1-4)MI^2 + \text{constante}$. Plus précisément, $MA^2 - 4MB^2 = -3MI^2 + (IA^2 - 4IB^2)$. En utilisant la relation de Leibniz, on trouve que l'ensemble est le cercle de diamètre $[AB]$ (cercle d'Apollonius). En effet, $\dfrac{MA}{MB} = 2$ donne $M$ sur le cercle de diamètre $[AB]$ avec $A$ et $B$ fixes.

\corrige{de l'exercice 14}
$MA = 3MB$ équivaut à $MA^2 = 9MB^2$, soit $MA^2 - 9MB^2 = 0$. Soit $I$ le barycentre de $(A,1)$ et $(B,-9)$. Alors $MA^2 - 9MB^2 = (1-9)MI^2 + (IA^2 - 9IB^2) = -8MI^2 + \text{constante}$. L'ensemble est un cercle de centre $I$. Calculons $I$ : $I$ est tel que $\overrightarrow{AI} = \dfrac{9}{8}\overrightarrow{AB}$. Avec $A(-1,0)$ et $B(3,0)$, $I$ a pour coordonnées $\left(-1 + \dfrac{9}{8}\times 4, 0\right) = \left(-1 + \dfrac{9}{2}, 0\right) = \left(\dfrac{7}{2}, 0\right)$. Le rayon se calcule en prenant un point particulier, par exemple $M$ sur la droite $(AB)$ vérifiant $MA = 3MB$, soit $M$ entre $A$ et $B$ ou à l'extérieur. L'ensemble est le cercle d'Apollonius.

\corrige{de l'exercice 15}
$\dfrac{MA}{MB} = \dfrac{MC}{MA}$ équivaut à $MA^2 = MB \cdot MC$. C'est une équation de cercle. En posant $M$ variable, on peut montrer que l'ensemble est le cercle de diamètre $[BC]$ si $A$ est sur ce cercle, mais ici c'est plus général. En utilisant les coordonnées, on peut montrer que c'est un cercle passant par $A$.

\corrige{du problème 1}
\begin{enumerate}[label=\arabic*)]
\item $G$ barycentre de $(A,1)$, $(B,2)$ et $(C,3)$ : $x_G = \dfrac{1\times 2 + 2\times 4 + 3\times (-2)}{6} = \dfrac{2 + 8 - 6}{6} = \dfrac{4}{6} = \dfrac{2}{3}$, $y_G = \dfrac{1\times 1 + 2\times 5 + 3\times 3}{6} = \dfrac{1 + 10 + 9}{6} = \dfrac{20}{6} = \dfrac{10}{3}$. Donc $G\left(\dfrac{2}{3},\dfrac{10}{3}\right)$.
\item $\overrightarrow{MA} + 2\overrightarrow{MB} + 3\overrightarrow{MC} = (1+2+3)\overrightarrow{MG} = 6\overrightarrow{MG}$.
\item $\|6\overrightarrow{MG}\| = 6\sqrt{2}$ donne $MG = \sqrt{2}$. L'ensemble est le cercle de centre $G$ et de rayon $\sqrt{2}$.
\item $MA^2 + 2MB^2 + 3MC^2 = 6MG^2 + GA^2 + 2GB^2 + 3GC^2$. Calculons $GA^2 = \left(2-\dfrac{2}{3}\right)^2 + \left(1-\dfrac{10}{3}\right)^2 = \left(\dfrac{4}{3}\right)^2 + \left(-\dfrac{7}{3}\right)^2 = \dfrac{16}{9} + \dfrac{49}{9} = \dfrac{65}{9}$. $GB^2 = \left(4-\dfrac{2}{3}\right)^2 + \left(5-\dfrac{10}{3}\right)^2 = \left(\dfrac{10}{3}\right)^2 + \left(\dfrac{5}{3}\right)^2 = \dfrac{100}{9} + \dfrac{25}{9} = \dfrac{125}{9}$. $GC^2 = \left(-2-\dfrac{2}{3}\right)^2 + \left(3-\dfrac{10}{3}\right)^2 = \left(-\dfrac{8}{3}\right)^2 + \left(-\dfrac{1}{3}\right)^2 = \dfrac{64}{9} + \dfrac{1}{9} = \dfrac{65}{9}$. Donc $GA^2 + 2GB^2 + 3GC^2 = \dfrac{65}{9} + 2\times\dfrac{125}{9} + 3\times\dfrac{65}{9} = \dfrac{65 + 250 + 195}{9} = \dfrac{510}{9} = \dfrac{170}{3}$. L'équation $6MG^2 + \dfrac{170}{3} = 100$ donne $6MG^2 = 100 - \dfrac{170}{3} = \dfrac{300 - 170}{3} = \dfrac{130}{3}$, soit $MG^2 = \dfrac{130}{18} = \dfrac{65}{9}$, $MG = \dfrac{\sqrt{65}}{3}$. L'ensemble est le cercle de centre $G$ et de rayon $\dfrac{\sqrt{65}}{3}$.
\end{enumerate}

\corrige{du problème 2}
\begin{enumerate}[label=\arabic*)]
\item $MA^2 + MB^2 = 2MI^2 + \dfrac{AB^2}{2} = 2MI^2 + 8$. L'équation $2MI^2 + 8 = 20$ donne $MI^2 = 6$, $MI = \sqrt{6}$. $E_1$ est le cercle de centre $I$ (milieu de $[AB]$) et de rayon $\sqrt{6}$.
\item $\dfrac{MA}{MB} = 2$ équivaut à $MA^2 = 4MB^2$, soit $MA^2 - 4MB^2 = 0$. Soit $J$ le barycentre de $(A,1)$ et $(B,-4)$. Alors $MA^2 - 4MB^2 = -3MJ^2 + (JA^2 - 4JB^2)$. $J$ est tel que $\overrightarrow{AJ} = \dfrac{4}{3}\overrightarrow{AB}$. Avec $AB=4$, $AJ = \dfrac{16}{3}$, $JB = \dfrac{4}{3}$. $JA^2 - 4JB^2 = \left(\dfrac{16}{3}\right)^2 - 4\left(\dfrac{4}{3}\right)^2 = \dfrac{256}{9} - 4\times\dfrac{16}{9} = \dfrac{256 - 64}{9} = \dfrac{192}{9} = \dfrac{64}{3}$. L'équation $-3MJ^2 + \dfrac{64}{3} = 0$ donne $MJ^2 = \dfrac{64}{9}$, $MJ = \dfrac{8}{3}$. $E_2$ est le cercle de centre $J$ et de rayon $\dfrac{8}{3}$.
\item $ABC$ rectangle en $A$, $AB=4$, $AC=3$, donc $BC=5$. Soit $G$ l'isobarycentre de $A$, $B$, $C$. $MA^2 + MB^2 + MC^2 = 3MG^2 + GA^2 + GB^2 + GC^2$. $GA^2 + GB^2 + GC^2 = \dfrac{AB^2 + AC^2 + BC^2}{3} = \dfrac{16 + 9 + 25}{3} = \dfrac{50}{3}$. L'équation $3MG^2 + \dfrac{50}{3} = 50$ donne $MG^2 = \dfrac{100}{9}$, $MG = \dfrac{10}{3}$. $E_3$ est le cercle de centre $G$ et de rayon $\dfrac{10}{3}$.
\item Pour montrer la concourance, il faut vérifier que les trois cercles ont un point commun. Par exemple, le point $M$ tel que $M$ soit le pied de la hauteur issue de $A$ dans le triangle $ABC$ pourrait convenir, mais il faut calculer. On peut montrer que le point $M$ situé à l'intersection de $E_1$ et $E_2$ vérifie aussi $E_3$ par un calcul.
\end{enumerate}

\corrige{du problème 3}
\begin{enumerate}[label=\arabic*)]
\item $\overrightarrow{AI} = \dfrac{1}{3}\overrightarrow{AB}$ donne $3\overrightarrow{AI} = \overrightarrow{AB} = \overrightarrow{AI} + \overrightarrow{IB}$, soit $2\overrightarrow{AI} = \overrightarrow{IB}$, donc $2\overrightarrow{IA} + \overrightarrow{IB} = \vec{0}$. $I$ est le barycentre de $(A,2)$ et $(B,1)$.
\item $\overrightarrow{BJ} = \dfrac{2}{3}\overrightarrow{BC}$ donne $3\overrightarrow{BJ} = 2\overrightarrow{BC} = 2(\overrightarrow{BJ} + \overrightarrow{JC})$, soit $3\overrightarrow{BJ} = 2\overrightarrow{BJ} + 2\overrightarrow{JC}$, donc $\overrightarrow{BJ} = 2\overrightarrow{JC}$, soit $\overrightarrow{JB} + 2\overrightarrow{JC} = \vec{0}$. $J$ est le barycentre de $(B,1)$ et $(C,2)$.
\item $\overrightarrow{CK} = \dfrac{3}{4}\overrightarrow{CA}$ donne $4\overrightarrow{CK} = 3\overrightarrow{CA} = 3(\overrightarrow{CK} + \overrightarrow{KA})$, soit $4\overrightarrow{CK} = 3\overrightarrow{CK} + 3\overrightarrow{KA}$, donc $\overrightarrow{CK} = 3\overrightarrow{KA}$, soit $\overrightarrow{KC} + 3\overrightarrow{KA} = \vec{0}$. $K$ est le barycentre de $(C,1)$ et $(A,3)$.
\item $G$ barycentre de $(A,2)$, $(B,1)$ et $(C,3)$. En utilisant l'associativité, $I$ barycentre de $(A,2)$ et $(B,1)$, donc $G$ est le barycentre de $(I,3)$ et $(C,3)$, donc $G$ est le milieu de $[IC]$. De même, $J$ barycentre de $(B,1)$ et $(C,2)$, donc $G$ est le barycentre de $(A,2)$ et $(J,3)$, donc $G$ est sur $(AJ)$. Ainsi $G$ est l'intersection de $(AI)$ et $(BJ)$.
\item $2MA^2 + MB^2 + 3MC^2 = 2GA^2 + GB^2 + 3GC^2$ équivaut à $2(MA^2 - GA^2) + (MB^2 - GB^2) + 3(MC^2 - GC^2) = 0$. En utilisant la relation de Leibniz, $2(\overrightarrow{MG} + \overrightarrow{GA})^2 + \dots$, on obtient $6MG^2 = 0$, soit $M = G$. L'ensemble est réduit au point $G$.
\end{enumerate}